\documentclass[10pt,a4paper]{amsart}
\usepackage[margin=19mm]{geometry}
\usepackage{amsmath,amssymb,mathtools}
\usepackage[scr=rsfs,cal=euler]{mathalfa}
\usepackage{xfrac}
\usepackage[unicode,hidelinks,bookmarks=false]{hyperref}
\newcommand{\ie}{i.e.\ }
\newcommand{\cf}{cf.\ }
\newcommand{\resp}{respectively\ }
\newcommand{\Subsection}{\subsection}
\providecommand{\nobreakdash}{\nobreak\hbox{-}}
\setlength{\emergencystretch}{2em}
\multlinegap0pt
\usepackage[shortlabels]{enumitem}
\usepackage{amssymb}
\newtheorem{assumption}{Assumption}
\newtheorem{lemma}{Lemma}
\title{Asymptotics for partition functions of random normal matrix models with singularities}
\author{Kohei Noda}
\date{Source: arXiv:2609.25307v1 (CC BY 4.0)}
\begin{document}
\maketitle

\noindent\emph{Source and licence.} Asymptotics for partition functions of random normal matrix models with singularities. Version arXiv:2609.25307v1 (CC BY 4.0), distributed by arXiv under the Creative Commons Attribution 4.0 International licence, http://creativecommons.org/licenses/by/4.0/, as recorded in the arXiv licence metadata for this exact version and shown on the abstract page https://arxiv.org/abs/2609.25307v1. Full licence notice: \texttt{licenses/arxiv-2609.25307-CC-BY-4.0.txt}.\par\smallskip

\noindent\textit{Editorial scope.} Counted unit: the article's lemma j+1 (source0.tex lines 3639--3649). The article's complete original proof of it is delivered with it (source0.tex lines 3651--3798, one \texttt{proof} environment of 148 lines). No mathematics is written, repaired or re-derived: the statement and the proof are the article's own lines, in the article's own notation, and no retained line is substituted or re-flowed.  Retained uncounted as the source-local context the proof needs: lines 177--186; lines 3220--3223.  Nothing was omitted from any retained span, so the package carries no omission record.  No retained line contains a citation, so the wrapper carries no bibliography and no bibliography entry.  No retained line contains a figure environment, an image-inclusion command or drawing code, so no image is delivered and the package declares no asset dependency.  Editorial formatting only, in addition: an AMS \texttt{amsart} preamble that restates the article's own declarations and macros used by the retained lines, namely the environments \texttt{assumption}, \texttt{lemma} on separate counters, together with the standard packages those lines need.  The retained environments print in this document as Assumption 1, Lemma 1 in order of appearance, whereas the article prints the same items with the numbering of its own full document.  The complete native source of the article as distributed in the arXiv e-print for this exact version is bundled unchanged as \texttt{sources/211288/source0.tex}.
\begin{assumption}\label{Assumption_Q}
The potential is radial, $Q(z)=q(|z|)$, and satisfies the following conditions.
\begin{enumerate}[label=\textup{(\arabic*)}]
\item\label{Assumption 1}
$\displaystyle \liminf_{|z|\to\infty}\frac{Q(z)}{2\log|z|}>1$.
\item\label{Assumption 2}
$Q$ is globally subharmonic, is $C^\infty$-smooth in a neighborhood of its droplet, and is strictly subharmonic there.
\item The droplet is a disk centered at the origin.
\end{enumerate}
\end{assumption}

\begin{equation}
\label{def of tau k ast epsilon inequality}
    \frac{\tau_{\rho_1}}{1-\epsilon}<\frac{\tau_{\rho_2}}{1+\epsilon},\qquad \frac{\tau_{\rho_2}}{1-\epsilon}<1. 
\end{equation}

\begin{lemma}\label{lemma: g1+leq j leq j+1}
Uniformly for $g_{1,+}+1\leq j\leq j_{1,+}$,
\begin{align}
\begin{split}
h_{n,j}^{(\mathrm{an})}
&=
\frac{2\rho_1e^{\mathsf{k}(\rho_1)}e^{-nV_{\tau}(\rho_1)}}{\eta_1 n}
\bigg[1-\frac{\mathfrak{a}_1}{n}+\frac{\mathfrak{a}_2}{n^2}-\frac{\mathfrak{a}_3}{n^3}+\mathcal{O}(\frac{1}{\eta_1^8n^4})\bigg].
\end{split}    
\end{align}
\end{lemma}

\begin{proof}
Split the integral into
\begin{align*}
h_{n,j}^{(\mathrm{an})}
&=\int_{0}^{\rho_1-\delta_n}2re^{\mathsf{k}(r)}e^{-nV_{\tau}(r)}\,dr
+\int_{\rho_1-\delta_n}^{\rho_1}2re^{\mathsf{k}(r)}e^{-nV_{\tau}(r)}\,dr
+\int_{\rho_2}^{+\infty}2re^{\mathsf{k}(r)}e^{-nV_{\tau}(r)}\,dr. 
\end{align*}
For $g_{1,+}+1\leq j\leq j_{1,+}$, we have
\[
V_{\tau}'(r)=\frac{1}{r}(rq'(r)-2\tau_j)\leq \frac{1}{r}\Bigl(rq'(r)-\frac{2}{n}\frac{n\tau_{\rho_1}}{1-\frac{M}{\sqrt{n}}}\Bigr).
\]
Because $Q$ is subharmonic in $\mathbb C$, the function $r\mapsto rq'(r)$ is increasing; moreover, $\rho_1q'(\rho_1)=2\tau_{\rho_1}$. Hence
\begin{equation}
\label{def of Vtau r leq rho1ast}
V_{\tau}'(r)
\leq \frac{1}{r}\Bigl(\rho_1q'(\rho_1)-\frac{2\tau_{\rho_1}}{1-\frac{M}{\sqrt n}}\Bigr)
=-\frac{2\tau_{\rho_1}}{r}\frac{\frac{M}{\sqrt n}}{1-\frac{M}{\sqrt n}},
\qquad r\leq \rho_1.
\end{equation}
Thus $V_\tau$ is decreasing on $(0,\rho_1]$, and $V_\tau(r)\geq V_\tau(\rho_1-\delta_n)$ for $r\leq\rho_1-\delta_n$. By the mean-value theorem,
\begin{equation}
    V_{\tau}(\rho_1-\delta_n)=V_{\tau}(\rho_1)-V_{\tau}'(r_1^{\ast})\delta_n,
\end{equation}
for some $r_{1}^{\ast}\in[\rho_1-\delta_n,\rho_1]$. 
Consequently, the first integral satisfies
\begin{align*}
\int_{0}^{\rho_1-\delta_n}2re^{\mathsf{k}(r)}e^{-nV_{\tau}(r)}\,dr
&\leq e^{-nV_{\tau}(\rho_1)}\mathcal O(e^{-cM^2})
\end{align*}
for some $c>0$ independent of $n$ and $j$. Thus this contribution is negligible uniformly in $j$.
 
We next show that the third integral is also negligible. Since
\begin{align*}
V_{\tau}(\rho_2)-V_{\tau}(\rho_1)
&=V_{\tau_{\rho_1}}(\rho_2)-V_{\tau_{\rho_1}}(\rho_1)+2(\tau_{\rho_1}-\tau)\log \frac{\rho_2}{\rho_1},
\end{align*}
we have 
\begin{equation}
\label{def of Vtau minus Vtau rho1ast}
V_{\tau}(r)-V_{\tau}(\rho_1) 
=
V_{\tau}(r)-V_{\tau}(\rho_2)
+
V_{\tau_{\rho_1}}(\rho_2)-V_{\tau_{\rho_1}}(\rho_1)+2(\tau_{\rho_1}-\tau)\log \frac{\rho_2}{\rho_1}.
\end{equation}
For $g_{1,+}+1\leq j\leq j_{1,+}$, we have $\tau_{\rho_1}-\tau>-\frac{\epsilon}{1-\epsilon}\tau_{\rho_1}$,
and hence
\begin{align}
    \begin{split}
\label{def of Vtau1ast rho2 - Vtau1ast rho1}
V_{\tau_{\rho_1}}(\rho_2)-V_{\tau_{\rho_1}}(\rho_1)+2(\tau_{\rho_1}-\tau)\log\frac{\rho_2}{\rho_1}
&\geq q(\rho_2)-q(\rho_1)-\frac{2\tau_{\rho_1}}{1-\epsilon}\log\frac{\rho_2}{\rho_1}
\\
&=
2\log\frac{\rho_2}{\rho_1}
\Bigl(\frac{q(\rho_2)-q(\rho_1)}{2\log\frac{\rho_2}{\rho_1}}-\frac{\tau_{\rho_1}}{1-\epsilon}\Bigr)
\\
&=2\log\frac{\rho_2}{\rho_1}
\Bigl(\sigma_{\star}-\frac{\tau_{\rho_1}}{1-\epsilon}\Bigr)
=:c_{\epsilon}>0
    \end{split}
\end{align}
for some $c_\epsilon>0$ independent of $n$ and $j$. Moreover, \eqref{def of tau k ast epsilon inequality} implies
\begin{align*}
 V_{\tau}'(r)&\geq \frac{1}{r}\Bigl(\rho_2q'(\rho_2)-\frac{2\tau_{\rho_1}}{1-\epsilon}\Bigr)
 \\
 &=
 \frac{1}{r}\Bigl(\frac{(1+\epsilon)\rho_2q'(\rho_2)}{1+\epsilon}-\frac{2\tau_{\rho_1}}{1-\epsilon}\Bigr)
 >\frac{1}{r}\Bigl(\frac{2(1+\epsilon)\tau_{\rho_1}}{1-\epsilon}-\frac{2\tau_{\rho_1}}{1-\epsilon}\Bigr)=\frac{1}{r}\frac{2\epsilon\tau_{\rho_1}}{1-\epsilon}>0,
\end{align*}
Thus $V_\tau$ is increasing on $[\rho_2,\infty)$ throughout the stated range of $j$, so $V_\tau(r)\geq V_\tau(\rho_2)$ for $r\geq\rho_2$, with equality only at $r=\rho_2$. Combining this observation with \eqref{def of Vtau minus Vtau rho1ast} and \eqref{def of Vtau1ast rho2 - Vtau1ast rho1} gives
\begin{equation}
\label{def of inequality Vtau r Vtau rho1ast}
    V_{\tau}(r)-V_{\tau}(\rho_1)>V_{\tau}(r)- V_{\tau}(\rho_2)+c_{\epsilon}.
\end{equation}
Therefore we obtain
\begin{align*}
    \int_{\rho_2}^{+\infty}2re^{\mathsf{k}(r)}e^{-nV_{\tau}(r)}\,dr
    &\leq 
    C_s e^{-nV_{\tau}(\rho_1)}
    e^{-nc_\epsilon}\int_{\rho_2}^{+\infty}r^{2\alpha+1}e^{-n(V_{\tau}(r)-V_{\tau}(\rho_2))}\,dr
    \\
    &\leq 
   C_s e^{-nV_{\tau}(\rho_1)}e^{-nc_{\epsilon}}
    \int_{\rho_2}^{+\infty}r^{2\alpha+1}e^{-p(V_{\tau}(r)-V_{\tau}(\rho_2))}\,dr
    =e^{-nV_{\tau}(\rho_1)}\cdot \mathcal{O}(e^{-nc_{\epsilon}'}),
\end{align*}
for $n\geq p$ and some $c_\epsilon'>0$.  Here $p$ is a fixed positive
integer chosen as follows.  Assumption~\ref{Assumption 1} gives
$q(r)\geq2(1+\varepsilon_0)\log r$ for all sufficiently large $r$ and
some $\varepsilon_0>0$.  Choose $p$ so that $p\varepsilon_0>\alpha+1$.
Since $0\leq\tau\leq1$, the last integrand is then bounded at infinity
by a constant times $r^{2\alpha+1-2p\varepsilon_0}$, uniformly in $j$.
On compact intervals uniform integrability follows from continuity.
The displayed integral is therefore uniformly bounded.

It remains to evaluate the second integral. Since $\rho_1q'(\rho_1)=2\tau_{\rho_1}$, for $g_{1,+}+1\leq j\leq j_{1,+}$ we have
\begin{equation}
\label{def of rho1 hard tau relationship}
-\frac{\epsilon\tau_{\rho_1}}{1-\epsilon}
\leq
\tau_{\rho_1}-\frac{\tau_{\rho_1}}{1-\epsilon}
\leq
\frac{1}{2}\rho_1q'(\rho_1)-\tau_j
\leq \tau_{\rho_1}-\frac{\tau_{\rho_1}}{1-\frac{M}{\sqrt{n}}}
=
-\frac{M}{\sqrt{n}}\frac{\tau_{\rho_1}}{1-\frac{M}{\sqrt{n}}}<0.
\end{equation}
Consequently, for $\eta_1\equiv\eta_1(j):=\frac{2}{\rho_1}(\tau-\tau_{\rho_1})>0$, one has $\eta_1^{-k}=\mathcal{O}(n^{k/2}M^{-k})$, uniformly in $j$. With the change of variables $u=n\eta_1(\rho_1-r)$, the endpoint Laplace expansion becomes
\begin{align*}
&\quad 2e^{\mathsf{k}(\rho_1)}e^{-nV_{\tau}(\rho_1)}\int_{\rho_1-\delta_n}^{\rho_1}re^{\mathsf{k}(r)-\mathsf{k}(\rho_1)}e^{-n(V_{\tau}(r)-V_{\tau}(\rho_1))}\,dr
\\
&=
\frac{2\rho_1e^{\mathsf{k}(\rho_1)}e^{-nV_{\tau}(\rho_1)}}{\eta_1 n}\int_0^{\eta_1n\delta_n}
e^{-u}
\Bigl(1-\frac{a_1(\eta_1^{-1}u)}{n}+\frac{a_2(\eta_1^{-1}u)}{n^2}-\frac{a_3(\eta_1^{-1}u)}{n^3}+\mathcal{O}(\frac{a_4(\eta_1^{-1}u)}{n^4})\Bigr)
\,du,
\end{align*}
where 
\begin{align*}
a_1(u)&:=\frac{1}{2}v_{2,1}u^2+\Bigl(\mathsf{k}'(\rho_1)+\frac{1}{\rho_1}\Bigr)u,
\\
a_2(u)&:=v_{2,1}^2\frac{u^4}{8}+\Bigl(\frac{3v_{2,1}}{\rho_1}+3v_{2,1}\mathsf{k}'(\rho_1)+v_{3,1}\Bigr)\frac{u^3}{6}
+\Bigl(\frac{2\mathsf{k}'(\rho_1)}{\rho_1}+\mathsf{k}'(\rho_1)^2+\mathsf{k}''(\rho_1)\Bigr)\frac{u^2}{2}
\\
a_3(u)&:=v_{2,1}^3\frac{u^6}{48}
+\Bigl(\frac{3v_{2,1}^2}{\rho_1}+3v_{2,1}^2\mathsf{k}'(\rho_1)+2v_{2,1}v_{3,1}\Bigr)\frac{u^5}{24}
\\
&\quad
+\Bigl(v_{4,1}+\frac{4v_{3,1}}{\rho_1}+4v_{3,1}\mathsf{k}'(\rho_1)+6\mathsf{k}''(\rho_1)v_{2,1}+6\mathsf{k}'(\rho_1)^2v_{2,1}+\frac{12v_{2,1}}{\rho_1}\mathsf{k}'(\rho_1)\Bigr)\frac{u^4}{24}
\\
&\quad
+\Bigl(\frac{3\mathsf{k}'(\rho_1)^2}{\rho_1}+\frac{3\mathsf{k}''(\rho_1)}{\rho_1}+\mathsf{k}'(\rho_1)^3+3\mathsf{k}'(\rho_1)\mathsf{k}''(\rho_1)+\mathsf{k}^{(3)}(\rho_1)\Bigr)
\frac{u^3}{6},
\end{align*}
and $a_4(u)$ is a polynomial in $u^8,u^7,u^6,u^5,u^4$. Termwise integration against $e^{-u}$ gives
\begin{align*}
&\quad 2e^{\mathsf{k}(\rho_1)}e^{-nV_{\tau}(\rho_1)}\int_{\rho_1-\delta_n}^{\rho_1}re^{\mathsf{k}(r)-\mathsf{k}(\rho_1)}e^{-n(V_{\tau}(r)-V_{\tau}(\rho_1))}\,dr
\\
&=
\frac{2\rho_1e^{\mathsf{k}(\rho_1)}e^{-nV_{\tau}(\rho_1)}}{\eta_1 n}
\bigg[1-\frac{\mathfrak{a}_1}{n}+\frac{\mathfrak{a}_2}{n^2}-\frac{\mathfrak{a}_3}{n^3}+\mathcal{O}(\frac{1}{\eta_1^8n^4})\bigg]. 
\end{align*}
By \eqref{def of rho1 hard tau relationship}, the remainder is $\mathcal{O}(\eta_1^{-8}n^{-4})=\mathcal{O}(M^{-8})$, uniformly over the full index range. This proves the lemma.


\end{proof}

\end{document}
